\documentclass[conference]{IEEEtran}
\IEEEoverridecommandlockouts
% The preceding line is only needed to identify funding in the first footnote. If that is unneeded, please comment it out.
\usepackage{cite}
\usepackage{amsmath,amssymb,amsfonts}
\usepackage{algorithmic}
\usepackage{graphicx}
\usepackage{textcomp}
\usepackage{xcolor}
\def\BibTeX{{\rm B\kern-.05em{\sc i\kern-.025em b}\kern-.08em
    T\kern-.1667em\lower.7ex\hbox{E}\kern-.125emX}}
\begin{document}

\title{Julia optimization analysis}

\author{\IEEEauthorblockN{1\textsuperscript{st} Rafał Pyłko}
\IEEEauthorblockA{\textit{Department of Electrical Engineering} \\
\textit{Warsaw University of Technology}\\
Warsaw, Poland \\
rafal.pylko.stud@pw.edu.pl}
\and
\IEEEauthorblockN{2\textsuperscript{nd} Michał Jagodziński}
\IEEEauthorblockA{\textit{Department of Electrical Engineering} \\
\textit{Warsaw University of Technology}\\
Warsaw, Poland \\
michal.jagodzinski2.stud@pw.edu.pl}
}

\maketitle

\begin{abstract}
Abstract about libraries comparison.
\end{abstract}

\begin{IEEEkeywords}
Julia, CNN, optimization, PyTorch, Neural Networks, SDG, Auto Differentiation
\end{IEEEkeywords}

\section{Introduction}
This document is a model and instructions for \LaTeX.
Please observe the conference page limits. 

\section{Optimization applied in our library}

\subsection{Types instability}

\subsection{im2col}

\subsection{float32}

\subsection{Weights initialization}

\subsection{In-place operations}

\section{Comparison with other libraries}

\subsection{Time complexity}

\subsection{Memory usage}

\section{Conclusion}
Julia is fast.

\begin{thebibliography}{00}
\bibitem{b1} G. Eason, B. Noble, and I. N. Sneddon, ``On certain integrals of Lipschitz-Hankel type involving products of Bessel functions,'' Phil. Trans. Roy. Soc. London, vol. A247, pp. 529--551, April 1955.
\end{thebibliography}

\end{document}
